\section{Consolidation: replaying the store into the weights}
\label{sec:sleep}

If the store is the state (\S\ref{sec:transplant}), it grows without bound; a memory
system needs a way to move old episodes out of the store and into the weights, and to
reclaim slots. We add a \emph{sleep} phase---an offline pass that replays store contents
and consolidates them---and characterise it (r26--r33). Two properties of the exact-token
key make this cheap.

\paragraph{The store replays itself, with no decoder.} A stored value does not decode
back to a token (a linear probe on values reads symbols at $0\%$), but a stored
\emph{key} does, exactly ($100\%$)---it is a token embedding. And because a value is
shifted by one, slot $i{+}1$'s key is slot $i$'s value token, so \emph{the sequence of
decoded keys reconstructs the written text} with no separately trained decoder (r26).
Replay is therefore free: the store is its own transcript.

\paragraph{Sleep consolidates and reclaims.} Replaying that transcript as a
language-model target moves the episode into the weights: recall from weights alone goes
$0\%\to\recall{100}$ and every one of the store's slots becomes reclaimable, with the
language-model loss \emph{improving} under rehearsal ($0.96\to0.83$) (r27). Rehearsing
the memory \emph{task} during sleep---not only the language stream---keeps the write path
calibrated so post-sleep new writes stay at \recall{100} (from $75\%$ without it) and
consolidation holds at the $8\times$ store ($3937$ slots) (r28). Sleep does age the key
space, though: consolidating shifts the embeddings a store's old keys were written
against, so an untouched old store decays ($88\%\to38\%$)---the key space cannot both
stay put for old stores and move for learning. Re-encoding surviving episodes from their
own decoded transcripts resolves this (r28).

\paragraph{Steady state needs bounded replay.} Across twelve wake/sleep cycles the loop
is stable---consolidation stays $88$--$100\%$, key drift is sub-linear and stays under
the destructive threshold, the language-model loss keeps improving---but knowledge is
not retained for free: without replay each episode is forgotten within about three
cycles. A fixed per-night replay budget fixes it, lifting twelve-user retention from
$0\%$ to a mean of $69\%$ at constant cost (r29).

\paragraph{Compression is free along one axis only.} Two ways to shrink an episode in
the store behave oppositely. Writing \emph{fewer slots} by a content-aware surprisal gate
is free---keeping the highest-surprisal quarter of tokens holds \recall{100}
($4\times$ fewer slots, baseline $99.7\%$), while a fixed stride at the same rate
collapses ($58\%$ / $0.5\%$) and an episode-end-only write is at chance (r30). Shrinking
the slot \emph{width} with a random projection, by contrast, fails at chance---it
destroys the exact-token match that the layer-0 key depends on (r31), the same reason a
contextualised key fails elsewhere. Treating episodes as first-class objects with
completeness-aware eviction beats ring eviction ($0.5\%\to\recall{90.6}$), given room for
one complete plus one replaceable episode (r31). Finally, compression interacts with
sleep: consolidation on a compressed store is unaffected, but the re-encode repair path
must carry the original written positions or it compresses twice and drops exactly the
fact tokens (the surprisal gate, after sleep, scores facts as low-surprisal because the
model now predicts them)---carrying positions restores it ($6.3\%\to75\%$), and
compression usefully removes the seed variance of full-rate writing ($97.3\%$ mean at an
$8$-point spread versus $71\%$ at $72$) (r32--r33).
